\subsection*{Luis Quintero: CAP/PACELC y comparación (3 minutos)}

\textbf{0:00--0:55.} No voy a etiquetar a Netflix entero como AP o EL. Analizo
una operación concreta: lecturas locales y replicación entre regiones de
EVCache para datos no críticos. Netflix permite que la réplica remota se atrase
sin detener las operaciones locales. Ante una partición mantiene disponibilidad
y tolera temporalmente valores distintos entre regiones.

\textbf{0:55--1:35.} Sin partición, la misma decisión favorece latencia: se lee
localmente y se replica de forma asíncrona. Para ese caso limitado hablamos de
AP/EL. No depende de gRPC. Otros datos de Netflix usan cuórum y consistencia de
lectura después de escritura, así que una sola etiqueta no sirve para toda la
plataforma.

\textbf{1:35--2:30.} Como sistema cercano comparo la infraestructura de video
de YouTube/Google. Netflix publica una arquitectura de microservicios, cachés y
contratos gRPC. Google documentó software distribuido coordinado con
aceleradores VCU para transcodificar video y reportó entre 20 y 33 veces más
eficiencia. No afirmamos que YouTube use aquí gRPC porque la fuente no lo dice.

\textbf{2:30--3:00.} RPC es el modelo general; gRPC es un framework con
contratos protobuf y soporte multilenguaje; RMI es una alternativa centrada en
Java. La conclusión es que la tecnología de comunicación y la política de
consistencia son decisiones distintas.
